\documentclass[11pt]{article}
\usepackage[utf8]{inputenc}
\usepackage[spanish,es-tabla]{babel}
\usepackage{amsmath,cancel}
\usepackage{amsfonts}
\usepackage{amssymb}
\usepackage{listings}
\usepackage{xcolor}
\usepackage{tabularx}
\usepackage{adjustbox}
\spanishdecimal{.}
\usepackage{graphicx}
\usepackage{float}
\usepackage[lmargin=2.5cm,rmargin=2.5cm,top=2.5cm,
bottom=2cm]{geometry}
\newcommand{\nombre}{Oscar}
\newcommand{\fecha}{3 de mayo del 2026}
\newcommand{\tituloPractica}{}
%%%%%%%%%%%%%%%%%%%%%%%%%%%%%%%%%%%%%%%%%%%%%%
\begin{document}
\begin{titlepage}
\begin{center}
\vspace*{2\baselineskip}
\hrule height 3pt
\vspace*{0.5\baselineskip}
{\Large \textbf{UNIVERSIDAD DE GUADALAJARA}}\\[0.1cm]
{\large \textbf{CENTRO UNIVERSITARIO DE CIENCIAS EXACTAS E INGENIER\'IAS}}
\vspace*{0.5\baselineskip}
\hrule
\vspace*{3\baselineskip}
\includegraphics[width=0.2\textwidth]{udg}
\vspace*{2\baselineskip}\\
{\large\textbf{ Algor\'itmos Bio-Inspirados \\ Msc. Rob\'otica e Inteligencia Artificial}
\vspace*{\baselineskip}\\
{\large Algoritmo 5 --- Particle Swarm Optimization N-Dimensional \vspace{1cm} \\
\begin{tabular}{p{4cm}p{8cm}}
\hline \hline
Nombre del alumno:  & Oscar Alberto Gallo Garc\'ia \\
C\'odigo del alumno:  & 218357704 \\
Profesor: & Dr. Carlos Alberto Lopez Franco \\
 \hline \hline
\end{tabular}
\vfill
\fecha
\end{center}
\newpage
\end{titlepage}

\tableofcontents
\newpage

\section{Resumen}

Se implement\'o el algoritmo de Optimizaci\'on por Enjambre de Part\'iculas (\textit{Particle Swarm Optimization}, PSO) en $n$ dimensiones, variante \textit{global best} (gbest). El enjambre est\'a compuesto por $n_s$ part\'iculas que se mueven en el espacio de b\'usqueda guiadas por su mejor posici\'on hist\'orica personal y por la mejor posici\'on encontrada por todo el enjambre. Se evaluaron doce casos: minimizaci\'on y maximizaci\'on en 2~dimensiones con gr\'aficas de superficie y contorno, y minimizaci\'on y maximizaci\'on en 10~dimensiones.

\section{Metodolog\'ia}

El PSO est\'a inspirado en el comportamiento colectivo de bandadas de p\'ajaros y cardumes de peces. Cada part\'icula representa una soluci\'on candidata en el espacio de b\'usqueda y ajusta su trayectoria en funci\'on de su propia experiencia y de la experiencia del enjambre.

\subsection{Posici\'on y velocidad}

La posici\'on de la part\'icula $i$ se actualiza a\~nadiendo su velocidad:
\[
\mathbf{x}_i(t+1) = \mathbf{x}_i(t) + \mathbf{v}_i(t+1)
\]

La velocidad incorpora tres componentes: inercia, cognitivo y social:
\[
v_{ij}(t+1) = w\,v_{ij}(t)
  + c_1\,r_{1j}(t)\bigl[y_{ij}(t) - x_{ij}(t)\bigr]
  + c_2\,r_{2j}(t)\bigl[\hat{y}_j(t) - x_{ij}(t)\bigr]
\]

donde $v_{ij}(t)$ es la velocidad de la part\'icula $i$ en la dimensi\'on $j$, $w$ es el peso de inercia, $c_1$ y $c_2$ son aceleraciones cognitiva y social respectivamente, y $r_{1j}(t),\,r_{2j}(t)\sim\mathcal{U}(0,1)$ aportan el componente estoc\'astico.

\subsection{Mejor personal}

El \textit{personal best} $\mathbf{y}_i$ es la mejor posici\'on visitada por la part\'icula $i$ desde la primera iteraci\'on. Para un problema de minimizaci\'on:
\[
\mathbf{y}_i(t+1) =
\begin{cases}
\mathbf{y}_i(t)   & \text{si } f\!\left(\mathbf{x}_i(t+1)\right) \geq f\!\left(\mathbf{y}_i(t)\right) \\
\mathbf{x}_i(t+1) & \text{si } f\!\left(\mathbf{x}_i(t+1)\right) < f\!\left(\mathbf{y}_i(t)\right)
\end{cases}
\]

\subsection{Mejor global}

El \textit{global best} $\hat{\mathbf{y}}(t)$ es la mejor posici\'on encontrada por cualquier part\'icula del enjambre:
\[
\hat{\mathbf{y}}(t) \in \{\mathbf{y}_0(t),\ldots,\mathbf{y}_{n_s}(t)\}
\;\Big|\;
f\!\left(\hat{\mathbf{y}}(t)\right) = \min\bigl\{f\!\left(\mathbf{y}_0(t)\right),\ldots,f\!\left(\mathbf{y}_{n_s}(t)\right)\bigr\}
\]

\subsection{Algoritmo gbest}

\begin{enumerate}
    \item Inicializar el enjambre: $\mathbf{x}_i \sim \mathcal{U}(\ell_{\min},\ell_{\max})$,\quad $\mathbf{v}_i \sim \mathcal{U}(-\Delta,\Delta)$,\quad $\mathbf{y}_i = \mathbf{x}_i$.
    \item Calcular $f(\mathbf{y}_i)$ para toda $i$ y determinar $\hat{\mathbf{y}}$.
    \item Repetir hasta la condici\'on de paro:
    \begin{enumerate}
        \item Para cada part\'icula $i$: actualizar $\mathbf{y}_i$ si $\mathbf{x}_i$ mejora; actualizar $\hat{\mathbf{y}}$ si $\mathbf{y}_i$ mejora.
        \item Para cada part\'icula $i$: actualizar $\mathbf{v}_i$ y $\mathbf{x}_i$; proyectar $\mathbf{x}_i$ al dominio $[\ell_{\min},\ell_{\max}]$.
    \end{enumerate}
    \item Retornar $\hat{\mathbf{y}}$.
\end{enumerate}

\subsection{Par\'ametros}

\begin{table}[H]
\centering
\begin{tabular}{lcc}
\hline\hline
Par\'ametro & 2D & 10D \\
\hline
$w$                       & 0.7  & 0.7  \\
$c_1$                     & 1.5  & 1.5  \\
$c_2$                     & 1.5  & 1.5  \\
$n_s$ (tama\~no enjambre) & 50   & 100  \\
Iteraciones m\'aximas     & 200  & 1000 \\
\hline\hline
\end{tabular}
\caption{Par\'ametros utilizados en todos los experimentos.}
\end{table}

\subsection{Funciones de prueba}

\begin{align}
f_{\text{esfera}}(\mathbf{x}) &= \sum_{j=1}^{n} x_j^2 \notag\\[4pt]
f_{\text{rastrigin}}(\mathbf{x}) &= 10n + \sum_{j=1}^{n}\bigl[x_j^2 - 10\cos(2\pi x_j)\bigr] \notag\\[4pt]
f_{\text{rosenbrock}}(\mathbf{x}) &= \sum_{j=1}^{n-1}\bigl[100(x_{j+1}-x_j^2)^2 + (1-x_j)^2\bigr] \notag\\[4pt]
f_{\text{ackley}}(\mathbf{x}) &= -20\exp\!\left(-0.2\sqrt{\tfrac{1}{n}\sum x_j^2}\right) - \exp\!\left(\tfrac{1}{n}\sum\cos(2\pi x_j)\right) + 20 + e \notag\\[4pt]
f_{\text{seno-coseno}}(\mathbf{x}) &= \sin(x_1)\cos(x_2) \notag\\[4pt]
f_{\text{easom}}(\mathbf{x}) &= \cos(x_1)\cos(x_2)\exp\!\bigl[-(x_1-\pi)^2-(x_2-\pi)^2\bigr] \notag
\end{align}

\section{Resultados}

\subsection{Minimizaci\'on 2D}

\subsubsection{Funci\'on Esfera}

M\'inimo global en $\mathbf{x}^* = (0,0)$ con $f(\mathbf{x}^*) = 0$. El PSO converge con precisi\'on num\'erica: mejor soluci\'on $(0,\,0)$, $f = 6.16\times10^{-25}$.

\begin{figure}[H]
\centering
\includegraphics[width=14cm]{figs/PSO_min_esfera.png}
\caption{Minimizaci\'on de la funci\'on Esfera. La estrella roja marca la mejor soluci\'on encontrada por el enjambre.}
\end{figure}

\subsubsection{Funci\'on de Rastrigin}

Funci\'on altamente multimodal con m\'inimo global en $\mathbf{x}^* = (0,0)$. El PSO escapa de los m\'inimos locales peri\'odicos y converge al \'optimo: mejor soluci\'on $(-0,\,0)$, $f = 0.0$.

\begin{figure}[H]
\centering
\includegraphics[width=14cm]{figs/PSO_min_rastrigin.png}
\caption{Minimizaci\'on de la funci\'on de Rastrigin. La topolog\'ia de valles peri\'odicos requiere que el enjambre mantenga diversidad durante la exploraci\'on.}
\end{figure}

\subsubsection{Funci\'on de Rosenbrock}

Valle curvo estrecho con m\'inimo global en $\mathbf{x}^* = (1,1)$. El PSO navega correctamente el valle: mejor soluci\'on $(1,\,1)$, $f = 4.40\times10^{-12}$.

\begin{figure}[H]
\centering
\includegraphics[width=14cm]{figs/PSO_min_rosenbrock.png}
\caption{Minimizaci\'on de la funci\'on de Rosenbrock. La superficie en forma de pl\'atano curvado exige que las part\'iculas sigan el valle hasta el m\'inimo.}
\end{figure}

\subsection{Maximizaci\'on 2D}

\subsubsection{Esfera Negativa}

Se maximiza $-f_{\text{esfera}}(\mathbf{x})$, cuyo m\'aximo global es en $(0,0)$. Mejor soluci\'on $(0,\,0)$, $f = -6.22\times10^{-26}$.

\begin{figure}[H]
\centering
\includegraphics[width=14cm]{figs/PSO_max_neg_esfera.png}
\caption{Maximizaci\'on de la esfera negativa. El enjambre converge al punto de m\'aximo de la superficie invertida.}
\end{figure}

\subsubsection{Funci\'on Seno-Coseno}

$f = \sin(x_1)\cos(x_2)$ tiene m\'aximo global de $1.0$ en $(−\pi/2,\, -\pi)$ y sim\'etrico en $(\pi/2,\, 0)$, $(\pi/2,\, 2\pi)$, etc. El enjambre encuentra: mejor soluci\'on $(-1.5708,\, -3.1416)$, $f = 1.0$.

\begin{figure}[H]
\centering
\includegraphics[width=14cm]{figs/PSO_max_seno_coseno.png}
\caption{Maximizaci\'on de $\sin(x_1)\cos(x_2)$. La funci\'on presenta m\'ultiples m\'aximos id\'enticos; el PSO localiza uno de ellos con valor exacto $f=1$.}
\end{figure}

\subsubsection{Funci\'on Easom}

M\'aximo global en $(\pi,\,\pi)$ con $f = 1.0$. La funci\'on es casi cero en casi todo el dominio, lo que dificulta la orientaci\'on del enjambre. Mejor soluci\'on $(3.1416,\, 3.1416)$, $f = 1.0$.

\begin{figure}[H]
\centering
\includegraphics[width=14cm]{figs/PSO_max_easom.png}
\caption{Maximizaci\'on de Easom. El pico estrecho en $(\pi,\pi)$ requiere que las part\'iculas exploren un dominio mayoritariamente plano antes de localizar el m\'aximo.}
\end{figure}

\subsection{Minimizaci\'on N-Dimensional (10D)}

\begin{table}[H]
\centering
\begin{tabular}{lp{8.5cm}r}
\hline\hline
Funci\'on & Mejor soluci\'on & $f(\hat{\mathbf{y}})$ \\
\hline
Esfera    & $[0,\;0,\;0,\;0,\;0,\;0,\;0,\;0,\;0,\;0]$                                       & $6.31\times10^{-69}$ \\[3pt]
Rastrigin & $[0,\;-0,\;0,\;0.995,\;-1.990,\;-0,\;0.995,\;-0,\;0.995,\;0]$                  & $6.96$ \\[3pt]
Ackley    & $[0,\;-0,\;-0,\;-0,\;0,\;-0,\;-0,\;0,\;-0,\;-0]$                               & $4.00\times10^{-15}$ \\
\hline\hline
\end{tabular}
\caption{Resultados de minimizaci\'on en 10 dimensiones. Esfera y Ackley convergen al \'optimo global con precisi\'on num\'erica. Rastrigin queda atrapada en un \'optimo local a $f=6.96$ (\'optimo global $= 0$).}
\end{table}

\subsection{Maximizaci\'on N-Dimensional (10D)}

Se maximizan las versiones negadas de las funciones anteriores.

\begin{table}[H]
\centering
\begin{tabular}{lp{8.5cm}r}
\hline\hline
Funci\'on & Mejor soluci\'on & $f(\hat{\mathbf{y}})$ \\
\hline
$-$Esfera    & $[0,\;-0,\;-0,\;-0,\;0,\;-0,\;0,\;0,\;-0,\;-0]$                             & $-1.08\times10^{-69}$ \\[3pt]
$-$Rastrigin & $[-0.995,\;-0,\;-0,\;-0,\;-1.990,\;0,\;-0,\;-0,\;0.995,\;0]$               & $-5.97$ \\[3pt]
$-$Ackley    & $[0,\;-0,\;0,\;-0,\;0,\;0,\;-0,\;0,\;-0,\;0]$                               & $-4.00\times10^{-15}$ \\
\hline\hline
\end{tabular}
\caption{Resultados de maximizaci\'on en 10 dimensiones. El comportamiento es sim\'etrico al de minimizaci\'on: $-$Esfera y $-$Ackley alcanzan el \'optimo global; $-$Rastrigin queda en un \'optimo local.}
\end{table}

\section{Conclusi\'on}

El PSO gbest encontr\'o el \'optimo global en funciones unimodales (Esfera, Rosenbrock, Ackley) tanto en 2D como en 10D con precisi\'on de m\'aquina, y localiz\'o los m\'aximos exactos de Seno-Coseno y Easom en 2D. En Rastrigin 10D el enjambre qued\'o en un \'optimo local ($f = 6.96$), lo que refleja la dificultad inherente de la funci\'on multimodal en alta dimensi\'on con los par\'ametros utilizados.\\\\
A diferencia del Recocido Simulado, el PSO explota informaci\'on colectiva: cada part\'icula no solo recuerda su propia mejor soluci\'on (\textit{personal best}) sino que se orienta por la mejor posici\'on del enjambre completo (\textit{global best}). Esto acelera la convergencia en espacios de alta dimensi\'on, pero puede reducir la diversidad y favorecer la convergencia prematura en funciones muy multimodales.\\\\
El peso de inercia $w$ controla el balance entre exploraci\'on y explotaci\'on: valores altos favorecen la b\'usqueda global y valores bajos refinan la soluci\'on. Los coeficientes $c_1$ y $c_2$ regulan qu\'e tanto conf\'ia cada part\'icula en su experiencia propia versus la del grupo; valores similares ($c_1 \approx c_2$) equilibran ambas influencias.

\end{document}
